\documentclass[../master.tex]{subfiles}

\begin{document}
\chapter{スタビライザー形式}
これまではShor符号の具体的な構成を通して、
量子誤り訂正符号の基本的なアイデアを紹介してきた。
このアプローチは、
情報を保護する量子状態を具体的にどのような形で符号化するかを先に定め、
その符号化を実現する回路を構成し、
その後に符号化された状態を壊さずに誤りを検出するための物理量を見つけるというものであった。
しかし、より高い誤り耐性などの性能を求めて複雑な符号を構築しようとする際、
この方法は非常に煩雑になり、いずれ限界を迎えることになる。
古典線形符号において符号は生成行列から作るものという視点だけではなく、
検査行列の核$(\ker H)$のなかの元という視点でも特徴づけることができる。
これと同様に量子誤り訂正においても、まずシンドローム測定のための物理量を先に与え、
その満たすべき条件から論理状態のなす空間を定めるという立場も取ることができる。
本章で扱うスタビライザー形式はこの考えを一般化したものであり、
可換な一般化パウリ演算子の集合を用いて符号空間をその共通固有空間として記述する
\cite{Gottesman1997-ep,Daniel-Gottesman1999-eh,Bullock2007-jl}。
この枠組みによって、
エラーの検出及び訂正に関わるシンドローム測定、
論理状態間を区別する論理演算子の構造が整理され、
さらに符号空間を可換ハミルトニアンの基底状態のなす空間としてみる見方も得られる。

\section{符号語から検査条件へ}
これまでの Shor 符号の議論では、まず論理状態の具体的な形を与え、
その符号化を実現する符号化する回路を構成し、
その後に誤りを検出するためのシンドローム測定を考えてきた。
このような構成は量子誤り訂正の基本的な仕組みを理解する上では有用である。
実際、どのような量子状態が保護されるべきか、
またどのような測定がその状態を壊さずに誤りの情報だけを読み出すのかを、
具体例を通して直接確かめることができるからである。

しかし、符号距離などの性能や実装との相性を考慮するためより一般の量子誤り訂正符号を構成しようとすると、
このように符号語や回路を先に与える立場は急速に煩雑になる。
例えばCSS符号では、シフトエラーと位相エラーの両方の検出を同時に行える条件は、
符号の検査行列の関係式であるCSS条件 $H_sH_p^\mathsf{T}=0$ を満たすことである。
これでは符号語から始めた場合、生成行列$G$に対応する符号化を考えて、
そこから誤りを検出するための$H_s,H_p$を読み取り、
それがCSS条件を満たすかどうかを確認するというようになる。
このように符号語から始めるアプローチの場合、
量子誤り訂正符号としてそもそも使えるかどうかを確認するのでさえ、非常に遠回りになっており、
そこから性能を考慮して符号を改良することもさらに難しくなってしまう。

ここで古典線形符号の場合を思い出そう。
古典線形符号では、符号語のなす空間(符号空間)は生成行列 $G$ の像$\Im G$ として与えることができる一方で、
検査行列 $H$ の核 $\ker H$ としても与えることができる。
すなわち、符号はメッセージからどのように符号語を作るかという立場から見ることも、
どのような検査条件を満たすものを符号語とみなすかという立場から見ることもできる。
後者の立場では、まず $H$ を与えることで符号語空間そのものが定まり、
その空間の基底を選ぶことによって対応する生成行列 $G$ を構成できる。
したがって、生成行列による記述と検査行列による記述は、
同じ符号空間の2つの異なる見方にほかならない。

量子誤り訂正符号においても同様の発想をとることができる。
すなわち、先に論理状態や符号化の回路を与えるのではなく、
まず誤りのシンドロームを与える可換な検査演算子を定め、
それらを満たす状態全体を符号空間とみなすのである。
この立場をとれば、シンドローム測定の回路は
既に与えられた量子状態を調べるための補助的な装置ではなく、
むしろ符号空間そのものを定義する基本データとして理解される。

\section{スタビライザー群と正規化群}
ではシンドローム測定で測定する物理量をどのように与えるかを考えよう。
$\mathbb{F}_p$におけるHadamardテストでは、
ユニタリ演算子の固有値が$\omega^k = \exp(\frac{2\pi i}{p}k)$であれば、
測定するquditは確実に状態$\ket{-k}$になることを見た。
このことは、誤りのシンドロームを測定するための物理量は、
誤りのシンドロームを表すユニタリ演算子の固有値が$\omega^k$であれば、
誤りの検出には都合が良いことを示唆している。
このような固有値を取るようなユニタリ演算子は一般化パウリ演算子の組み合わせで構成される。
そのような演算子の集合を考える必要があるのとわかる。

また、誤りのシンドロームを測定するための物理量は測定する順番によって結果が変わらないようにする必要がある。
すなわち、測定する物理量は互いに可換でなければならない。
\begin{equation}\begin{aligned}[b]
    \mathcal{S}^{(1)}\mathcal{S}^{(2)}\ket{\psi} &= \mathcal{S}^{(2)}\mathcal{S}^{(1)}\ket{\psi}  &
    \Rightarrow \quad \mathcal{S}^{(1)}\mathcal{S}^{(2)}\mathcal{S}^{-1 (2)} \mathcal{S}^{-1 (1)}&=1
\end{aligned}\end{equation}
この条件より、測定する演算子の集合には演算が定義されていて、単位元や逆元を必要とする。
なので、測定する演算子の候補の集合は群をなす必要があり、これは一般化パウリ群として定義される。

\begin{NoteBox}[title={一般化パウリ群}]{}
    一般化パウリ群とは、$n$個のquditに作用する一般化パウリ演算子全体の集合である。
    \begin{equation}\begin{aligned}[b]
        \mathcal{P}_n = \{ \omega^k \mathcal{X}^{\mathbf{a}} \mathcal{Z}^{\mathbf{b}} \mid k \in \mathbb{Z}_p, \mathbf{a},\mathbf{b} \in \mathbb{F}_p^n \}
    \end{aligned}\end{equation}
\end{NoteBox}

そのうち、測定する演算子の集合は互いに可換でなければならないので、可換な部分群を考える必要がある。
\begin{NoteBox}[title={スタビライザー群}]{}
    スタビライザー群とは次の2つの条件を満たす一般化パウリ群$\mathcal{P}_n$の部分群である。
    \begin{itemize}
        \item $\mathcal{S} \subset \mathcal{P}_n$は可換な部分群である。
        \item $\omega I \notin \mathcal{S}$。
    \end{itemize}
\end{NoteBox}

スタビライザー群$\mathcal{S}$を定めると、$\mathcal{S}$の共通固有空間が符号空間となる。
\begin{NoteBox}[title={スタビライザー符号}]{}
    スタビライザー符号とは、スタビライザー群$\mathcal{S}$の共通固有空間を符号空間とする量子誤り訂正符号である。
    \begin{equation}\begin{aligned}[b]
        \mathcal{C} = \{ \ket{\psi} \in (\mathbb{C}^p)^{\otimes n} \mid S\ket{\psi} = \ket{\psi}, \forall S \in \mathcal{S} \}
    \end{aligned}\end{equation}
\end{NoteBox}


\end{document}
